%\section{Top quark mass scheme and scale uncertainties at NLO}
%\subsection{Top-quark mass effects}
\subsection{Top quark mass scheme and scale uncertainties at NLO}
\label{sec:NLO-uncertainties}

%\textit{Authors:} J. Baglio, F. Campanario, S. Glaus, M. M. Muhlleitner, J. Ronca, M. Spira, J. Streicher (2003.03227 + 1811.05692 + 2008.11626) \\[0.3cm]

\textbf{Julien Baglio, Francisco Campanario, Seraina Glaus, Milada Margarete Mühl\-leitner, Jonathan Ronca, Michael Spira, Juraj Streicher~\cite{Baglio:2018lrj,Baglio:2020ini,Baglio:2020wgt}.}


%\APtodo{Please decide if you would like to describe the NLO calculation in a preceding section (see~\ref{sec:nlo2})}
%        ==============================================
\noindent The analysis of the top quark mass scheme and scale uncertainties was first performed in Refs.~\cite{Baglio:2018lrj,Baglio:2020wgt,Baglio:2020ini}, including the full NLO QCD calculation (see Section~\ref{sec:nlo2}). The top quark mass scheme and
scale uncertainties drop by roughly a factor of two from LO to NLO. At
LO, we obtain the uncertainties ($Q$ denotes the invariant Higgs boson pair mass $m_{hh}$)
\begin{eqnarray}
\frac{d\sigma(gg\to hh)}{dQ}\Big|_{Q=300~{\rm GeV}} & = &
0.01656^{+62\%}_{-2.4\%}\, \mathrm{fb/GeV},\nonumber\\
\frac{d\sigma(gg\to hh)}{dQ}\Big|_{Q=400~{\rm GeV}} & = &
0.09391^{+0\%}_{-20\%}\, \mathrm{fb/GeV},\nonumber\\
\frac{d\sigma(gg\to hh)}{dQ}\Big|_{Q=600~{\rm GeV}} & = &
0.02132^{+0\%}_{-48\%}\, \mathrm{fb/GeV},\nonumber\\
\frac{d\sigma(gg\to hh)}{dQ}\Big|_{Q=1200~{\rm GeV}} & = &
0.0003223^{+0\%}_{-56\%}\, \mathrm{fb/GeV}
\end{eqnarray}
where the full spread of the cross sections is taken over predictions obtained with the top quark pole mass
(central values) and with the $\overline{\rm MS}$ mass. In the latter case, the scale
$\mu_t$ of $\overline{m}_t(\mu_t)$ is either set to $\overline{m}_t$ or
varied in the range between $Q/4$ and $Q$, as
in the single (off-shell) Higgs case considered earlier. The final NLO
results read \cite{Baglio:2018lrj,Baglio:2020wgt,Baglio:2020ini}
\begin{eqnarray}
\frac{d\sigma(gg\to hh)}{dQ}\Big|_{Q=300~{\rm GeV}} & = &
0.02978(7)^{+6\%}_{-34\%}\, \mathrm{fb/GeV},\nonumber\\
\frac{d\sigma(gg\to hh)}{dQ}\Big|_{Q=400~{\rm GeV}} & = &
0.1609(7)^{+0\%}_{-13\%}\, \mathrm{fb/GeV},\nonumber\\
\frac{d\sigma(gg\to hh)}{dQ}\Big|_{Q=600~{\rm GeV}} & = &
0.03204(9)^{+0\%}_{-30\%}\, \mathrm{fb/GeV},\nonumber\\
\frac{d\sigma(gg\to hh)}{dQ}\Big|_{Q=1200~{\rm GeV}} & = &
0.000435(6)^{+0\%}_{-35\%}\, \mathrm{fb/GeV}
\end{eqnarray}
Since these uncertainties are similar in size to the renormalization
and factorization scale dependencies at NLO, they constitute an important
contribution to the total theoretical uncertainties.

The amplitude may be written as the sum of
two form factors, $F_1$ and $F_2$, describing the scattering of incoming gluons with the
same helicity and opposite helicities, respectively. The contribution of the box
diagrams to the two form factors dominates at high energy.
Expanding the LO and NLO results of Ref.~\cite{Davies:2018qvx} around
large invariant Higgs boson pair mass, $s$, we have, in the on-shell scheme and
in the notation of Ref.~\cite{Davies:2018qvx},
\begin{eqnarray}
F_\mathrm{box,i} & = & F_\mathrm{box,i}^{(0)} +
\frac{\alpha_s(\mu_R)}{\pi} F_\mathrm{box,i}^{(1)} \qquad (i=1,2) \nonumber \\
F_\mathrm{box,i}^{(0)} & = & \frac{m_t^2}{s} c_{0,i} +
\mathcal{O}\left(\frac{1}{s^2}\right)  \nonumber \\
F_\mathrm{box,i}^{(1)} & = & 2 F_\mathrm{box,i}^{(0)} \log\frac{m_t^2}{s}
+ \frac{m_t^2}{s}  c_{1,i} +
\mathcal{O}\left(\frac{1}{s^2}\right)
\label{eq:hhexp}
\end{eqnarray}
where the coefficients $c_{0,i}$ and $c_{1,i}$ do not depend on the
top quark mass. The overall factor of $m_t^2$ for the box contribution
originates entirely from the Yukawa couplings, and examining $F_\mathrm{box,i}^{(0)}$,
we find that the form factors are independent of the propagator top quark mass in the high-energy limit.
Therefore, for a fixed Yukawa coupling, we expect the results in different schemes to asymptote.
Transforming the top quark pole mass $m_t$ into the $\overline{\rm MS}$ mass
$\overline{m}_t(\mu_t)$, the explicit expressions above are modified to
\cite{Baglio:2020ini}
\begin{eqnarray}
F_\mathrm{box,i}^{(0)} & = & \frac{\overline{m}_t^2(\mu_t)}{s} c_{0,i}
+ \mathcal{O}\left(\frac{1}{s^2}\right) \qquad (i=1,2) \nonumber \\
F_\mathrm{box,i}^{(1)} & = & 2 F_\mathrm{box,i}^{(0)}
\left[\log\frac{\mu_t^2}{s} + \frac{4}{3}\right] +
\frac{\overline{m}_t^2(\mu_t)}{s} c_{1,i} +
\mathcal{O}\left(\frac{1}{s^2}\right)
\label{eq:NLO-MSbar}
\end{eqnarray}
This emphasizes that, in order to absorb the large logarithmic terms $\log\frac{m_t^2}{s}$,
the choice $\mu_t=\kappa\sqrt{s}$ is preferred\footnote{A similar dynamical scale choice, e.g.~the transverse mass $\mu_t = \kappa M_T$ with $M_T = tu/s$ in the high-energy limit, would be equally favoured.} when the coefficient $\kappa$ is not too far from unity. However, the central recommended values of the total and differential cross sections use the top quark pole mass, so that this defines the reference prediction for the estimate of the uncertainties.

%\section{Expansion methods}

%\textit{Authors:} Roberto Bonciani, Giuseppe Degrassi, Pier Paolo Giardino, Ramona Gröber (1806.11564) + Emanuele Bagnaschi (2309.1052) + Davies et. al (1801.09696, 2302.01356 + 2506.04323)
%\APtodo{Decide what to do with the combination of the high-energy expansion with the NLO}
%        =================
%%An alternative approach to the fully numerical evaluation of the virtual corrections are expansions. In particular an expansion in small $p_T$ as originally proposed in \cite{Bonciani:2018omm} can cover more than 95\% of the relevant phase space for SM Higgs boson pair production. This expansion assumes that $p_T^2$ and $m_h^2$ are much smaller than $4 m_t^2$ allowing to Taylor expand in these quantities, while keeping the partonic centre-of-mass energy $\hat{s}$ arbitrary. This reduces the number of scales in the two-loop integrals to one which allows for the analytical determination. In Ref.~\cite{Bonciani:2018omm} all the two-loop integrals analytically in terms of generalised harmonic polylogarithms, but for two elliptic integrals that were determined in terms of series expansions around singular points \cite{Bonciani:2018uvv} following the approach of \cite{Pozzorini:2005ff, Aglietti:2007as}.

%In the high-energy limit the described expansion is no longer valid as large $p_T$ values can be achieved. For this very reason, this has to be combined with a complementary expansion as provided in \cite{Davies:2018qvx}. This reference expands the amplitudes in the high energy limit namely for $\hat{s}, \hat{t}, \hat{u} \gg m_t^2$. The expansion allows to express the integrals all analytically in terms of harmonic polylogarithms \cite{Davies:2018ood}.

%The two expansion have been combined by means of Pad\'e approximations in \cite{Bellafronte:2022jmo} showing that a very reliable result on the level of $< 1\%$ difference with the numerical grid of \cite{Davies:2019dfy} are obtained.

%This approach has been used for the POWHEG implementation of Higgs boson pair production at NLO QCD presented in \cite{Bagnaschi:2023rbx}. Due to the flexibility of the analytic approach the top mass renormalisation scheme uncertainty can be computed fully differentially, showing that in dependence of the kinematic variable considered it can be even larger in certain part of the phase space as for the inclusive cross section. The code so far allows to vary the top Yukawa coupling and the trilinear Higgs boson self-coupling but further extensions to beyond the SM scenarios can be easily achieved due to the flexibility of the approach.

%A similar approach has been followed for the public code \texttt{ggxy} \cite{Davies:2025qjr} relying on a deeper expansion in small Mandelstam $\hat{t}$ combined with a high energy expansion \cite{Davies:2023vmj}. The two-loop integrals have been obtained semi-analytically in terms of series expansions in one variable \cite{Fael:2021kyg}. In this approach, also first partial results for the virtual corrections at NNLO QCD have been obtained in \cite{Davies:2023obx, Davies:2024znp, Davies:2025ghl}

%Other approaches for instance include a combination of large mass expansion with a threshold expansion combined via Pad\'e approximants \cite{Grober:2017uho}, or an expansion in the Higgs mass $m_h$ with numerical evaluation of the two-loop integrals \cite{Xu:2018eos}.

%Finally, we want to emphasise that if one were to associate an uncertainty to the analytic approach with respect to the numeric one, it would be below 1\% and hence currently completely negligible with respect to other theory uncertainties.
